\sectionkicker{Source-backed technical notes}
\subsection*{モデル性能を実用能力に変えるServingとData: Technical Notes}
本欄は記事本文で圧縮した一次資料上の情報を、比較・再検証しやすい形で整理したものである。時系列、確認済みの主張、留意点、一次資料URLを掲載する。
\smallskip
\noindent{\small\textit{Technical Notesでは、一次資料から確認した公開・提供・研究上の事実を記載する。数値・能力評価や提供元・著者による比較表現は、独立再現済みの結果ではなく帰属claimとして扱う。}}\par
\medskip
\subsection*{Theme at a glance}
\begin{center}
\footnotesize
\begin{tabularx}{\linewidth}{@{}p{0.20\linewidth}p{0.10\linewidth}p{0.14\linewidth}X@{}}
\toprule
資料 & 位置づけ & 種別 & 時系列 \\
\midrule
NVIDIA NIM & 主要資料 & Framework & 2024-03-18             \\
FineWeb & 主要資料 & 論文 & 2024-06-25             \\
Infini-attention & 補足資料 & 論文 & 2024-04-10             \\
Meta Llama 3 & 補足資料 & オープンウェイト & 2024-04-18             \\
OLMo / OLMo 1.7 & 補足資料 & オープンウェイト & 2024-02-01, 2024-04-17             \\
\bottomrule
\end{tabularx}
\end{center}
\normalsize

\Needspace{0.38\textheight}
\begin{technicalnote}{NVIDIA NIM}{主要資料}
\begingroup
% reader-facing Technical Notes break policy
\widowpenalty=10000
\clubpenalty=10000
\displaywidowpenalty=10000
\begin{tabularx}{\linewidth}{@{}>{\bfseries}p{0.22\linewidth}X@{}}
組織 & NVIDIA \\
種別 & Framework \\
時系列 & 2024-03-18 \\
\end{tabularx}
\smallskip
\begin{minipage}{\linewidth}
% reader-facing Technical Notes event-only fact/source tail group
{\bfseries 一次資料から整理したtechnical points}
\begin{itemize}[leftmargin=1.5em,itemsep=0.35em]
\item \textbf{一次情報で確認できる事実}: NVIDIAの一次資料により、NVIDIA NIMについて2024-03-18に製品発表を確認した。 NVIDIA NIMはGTC 2024でinference microservicesとして提示され、NVIDIAのaccelerated-computing librariesとgenerative-AI modelsをパッケージ化し、industry-standard APIsで接続できる配備単位として説明された。NVIDIAはCUDA installed base上での利用、GPU世代に応じた再最適化、security vulnerability scanningを運用面の特徴として挙げている。
\end{itemize}
\begin{minipage}{\linewidth}
% reader-facing Technical Notes generic boundary/source tail group
\begin{samepage}
{\bfseries 一次資料}
\begingroup\sloppy
\Urlmuskip=0mu plus 2mu\relax
\begin{itemize}[leftmargin=1.5em,itemsep=0.25em]
\item {\scriptsize\url{https://blogs.nvidia.com/blog/2024-gtc-keynote/}}
\end{itemize}
\endgroup
\end{samepage}
\end{minipage}
\end{minipage}
\endgroup
\end{technicalnote}
\medskip

\Needspace{0.38\textheight}
\begin{technicalnote}{FineWeb}{主要資料}
\begingroup
% reader-facing Technical Notes break policy
\widowpenalty=10000
\clubpenalty=10000
\displaywidowpenalty=10000
\begin{tabularx}{\linewidth}{@{}>{\bfseries}p{0.22\linewidth}X@{}}
組織 & Hugging Face \\
種別 & 論文 \\
時系列 & 2024-06-25 \\
\end{tabularx}
\smallskip
\begin{minipage}{\linewidth}
% reader-facing Technical Notes event-only fact/source tail group
{\bfseries 一次資料から整理したtechnical points}
\begin{itemize}[leftmargin=1.5em,itemsep=0.35em]
\item \textbf{一次情報で確認できる事実}: Hugging Faceの一次資料により、FineWebについて2024-06-25にPaper First Submissionを確認した。 FineWebはCommon Crawl由来の大規模pretraining corpusを、URL/content filtering、language/quality filtering、repetition removal、deduplication等のdata-processing choicesを実験で比較しながら構築したopen datasetである。論文は約15T token規模のFineWebを公開し、各filtering choiceを小規模model training/evaluationで検証するdata-centricな設計を採る。下流性能値は著者報告として扱う。
\end{itemize}
\begin{minipage}{\linewidth}
% reader-facing Technical Notes generic boundary/source tail group
\begin{samepage}
{\bfseries 一次資料}
\begingroup\sloppy
\Urlmuskip=0mu plus 2mu\relax
\begin{itemize}[leftmargin=1.5em,itemsep=0.25em]
\item {\scriptsize\url{http://arxiv.org/abs/2406.17557v2}}
\end{itemize}
\endgroup
\end{samepage}
\end{minipage}
\end{minipage}
\endgroup
\end{technicalnote}
\medskip

\Needspace{0.38\textheight}
\begin{technicalnote}{Infini-attention}{補足資料}
\begingroup
% reader-facing Technical Notes break policy
\widowpenalty=10000
\clubpenalty=10000
\displaywidowpenalty=10000
\begin{tabularx}{\linewidth}{@{}>{\bfseries}p{0.22\linewidth}X@{}}
組織 & Google \\
種別 & 論文 \\
時系列 & 2024-04-10 \\
\end{tabularx}
\smallskip
\begin{minipage}{\linewidth}
% reader-facing Technical Notes event-only fact/source tail group
{\bfseries 一次資料から整理したtechnical points}
\begin{itemize}[leftmargin=1.5em,itemsep=0.35em]
\item \textbf{一次情報で確認できる事実}: Googleの一次資料により、Infini-attentionについて2024-04-10にPaper First Submissionを確認した。 Infini-attentionは、通常のattention blockへcompressive memoryを組み込み、masked local attentionとlong-term linear attentionを単一Transformer block内で扱う構成である。著者らはbounded memory/computationを保った長文脈処理とstreaming inferenceを目的として提示しており、長文脈タスク上の性能値は著者報告として扱う。
\end{itemize}
\begin{minipage}{\linewidth}
% reader-facing Technical Notes generic boundary/source tail group
\begin{samepage}
{\bfseries 一次資料}
\begingroup\sloppy
\Urlmuskip=0mu plus 2mu\relax
\begin{itemize}[leftmargin=1.5em,itemsep=0.25em]
\item {\scriptsize\url{http://arxiv.org/abs/2404.07143v2}}
\end{itemize}
\endgroup
\end{samepage}
\end{minipage}
\end{minipage}
\endgroup
\end{technicalnote}
\medskip

\Needspace{0.38\textheight}
\begin{technicalnote}{Meta Llama 3}{補足資料}
\begingroup
% reader-facing Technical Notes break policy
\widowpenalty=10000
\clubpenalty=10000
\displaywidowpenalty=10000
\begin{tabularx}{\linewidth}{@{}>{\bfseries}p{0.22\linewidth}X@{}}
組織 & Meta \\
種別 & オープンウェイト \\
時系列 & 2024-04-18 \\
\end{tabularx}
\smallskip
\begin{minipage}{\linewidth}
% reader-facing Technical Notes event-only fact/source tail group
{\bfseries 一次資料から整理したtechnical points}
\begin{itemize}[leftmargin=1.5em,itemsep=0.35em]
\item \textbf{一次情報で確認できる事実}: Metaの一次資料により、Meta Llama 3について2024-04-18にモデル公開を確認した。 Meta Llama 3の2024年4月公開分は8B/70Bのpretrained版とinstruction-fine-tuned版で、decoder-only Transformerを基盤に128K-token vocabularyのtokenizer、両サイズでgrouped query attention (GQA)、8,192-token sequence trainingを採用した。Metaはpretraining corpusを公開情報由来15T token超と説明しており、benchmark性能はMetaによる評価として扱う。
\end{itemize}
\begin{minipage}{\linewidth}
% reader-facing Technical Notes generic boundary/source tail group
\begin{samepage}
{\bfseries 一次資料}
\begingroup\sloppy
\Urlmuskip=0mu plus 2mu\relax
\begin{itemize}[leftmargin=1.5em,itemsep=0.25em]
\item {\scriptsize\url{https://ai.meta.com/blog/meta-llama-3/}}
\end{itemize}
\endgroup
\end{samepage}
\end{minipage}
\end{minipage}
\endgroup
\end{technicalnote}
\medskip

\Needspace{0.38\textheight}
\begin{technicalnote}{OLMo / OLMo 1.7}{補足資料}
\begingroup
% reader-facing Technical Notes break policy
\widowpenalty=10000
\clubpenalty=10000
\displaywidowpenalty=10000
\begin{tabularx}{\linewidth}{@{}>{\bfseries}p{0.22\linewidth}X@{}}
組織 & Ai2 \\
種別 & オープンウェイト \\
時系列 & 2024-02-01, 2024-04-17 \\
\end{tabularx}
\smallskip
\begin{minipage}{\linewidth}
% reader-facing Technical Notes event-only fact/source tail group
{\bfseries 一次資料から整理したtechnical points}
\begin{itemize}[leftmargin=1.5em,itemsep=0.35em]
\item \textbf{一次情報で確認できる事実}: Ai2の一次資料により、OLMo / OLMo 1.7について2024-02-01にモデル公開、2024-04-17にモデル更新を確認した。 対象event近傍の一次資料から 7B parameter scale を確認できる。
\end{itemize}
\begin{minipage}{\linewidth}
% reader-facing Technical Notes generic boundary/source tail group
\begin{samepage}
{\bfseries 一次資料}
\begingroup\sloppy
\Urlmuskip=0mu plus 2mu\relax
\begin{itemize}[leftmargin=1.5em,itemsep=0.25em]
\item {\scriptsize\url{https://allenai.org/blog/olmo-1-7-7b-a-24-point-improvement-on-mmlu-92b43f7d269d}}
\item {\scriptsize\url{https://allenai.org/blog/olmo-open-language-model-87ccfc95f580}}
\end{itemize}
\endgroup
\end{samepage}
\end{minipage}
\end{minipage}
\endgroup
\end{technicalnote}
\medskip
